\begin{abstract}
Telemetry reduction can remove the observations needed to detect or investigate a
fault. We study causal retention under an explicit serialized-byte budget and
separate capture decisions from assessment of decoded packets. A bounded recorder
combines pre-trigger context, post-trigger capture, and detector-independent
background sampling. Comparisons use native observations, source-fixed event
denominators, and training-calibrated receivers that remain frozen across policies.
On 3,248,700 Mission1 validation observations containing 11 eligible events, the
recorder preserves three-component PCA evidence for 11/11 events using 6.30 MiB,
compared with 5/11 using 25.04 MiB for uniform sampling at the same 1,024-byte
hourly cap. Removing pre-trigger capture leaves the event count unchanged but
reduces pre-event context coverage from 64.3\% to 1.8\%. Seeded controls,
component ablations, and a 132-arm capture/receiver comparison examine the scope
of these results. A separate AR(1) receiver exposes a temporal limitation:
with the trigger disabled, singleton background sampling preserves PCA evidence
for 4/11 events but temporal evidence for 0/11. An adjacent-pair background control,
developed after inspecting this failure, recovers temporal evidence for 2/11.
Illustrative finite-buffer delivery scenarios further distinguish captured from
delivered evidence. The results show why event counts, context, receiver input
requirements, and actual expenditure must be evaluated together. They remain
development evidence on one mission; the exploratory pair control, independent-
mission generalization, and operational diagnostic usefulness require further
validation.
\end{abstract}
